\chapter{Bounds, Minimum and Maximum}

\section{Finding Tight Intervals}
Recall the definition of $\min$ and $\max$ functions from the chapter \textbf{Range Manipulation}'s \textbf{Value Clamping} section.

We can generalize the $\min$ and $\max$ function to include more values :3

For $\min$:
\begin{center}
        $\min(a_1, a_2, a_3, ..., a_n) = \min(a_1, \min(a_2, a_3, ..., a_n))$
\end{center}

And for $\max$:
\begin{center}
        $\max(a_1, a_2, a_3, ..., a_n) = \max(a_1, \max(a_2, a_3, ..., a_n))$
\end{center}

For example, $\min(1, 4, -4) = \min(1, \min(4, -4))$. Here we can evaluate from the inside to the outside.
So $\min(4, -4) = -4$ and finally $\min(1, -4) = -4 = \min(1, 4, -4)$. Same should apply to the generalized $\max$ function!

In other words:
\begin{itemize}
    \item $\min$: Returns the minimum among all the values $a_i$.
    \item $\max$: Returns the maximum among all the values $a_i$.
\end{itemize}

By using this generalized $\min$ and $\max$ functions, we can actually find the bounding range of a set of values,
for example, if we have a set of values denoted by $a_i$, and there are $n$ of them (i.e. the first element is $a_1$ and last is $a_n$),
then, the interval that \textbf{tightly} includes all values is defined as:
\begin{center}
    $[\min(a_1, a_2, a_3, ..., a_n), \max(a_1, a_2, a_3, ..., a_n)]$
\end{center}

For example, if we have values $1, 4, 5, -2, 3$, we can find it's tight bound as:
\begin{center}
    $[\min(1, 4, 5, -2, 3), \max(1, 4, 5, -2, 3)] = [-2, 5]$
\end{center}

Here's how it looks like when plotted in the real number line:

TODO: an image here demonstrating the example :3



\section{Finding Bounds in Higher Dimensions}
This can be very useful if we extend this to two or more dimensions, for instance, we can find the bounding box in two dimensions
given a bunch of points by finding the $\min$ and $\max$ of all points:

TODO: axis-aligned bounding boxes etc.





